\begin{proposition}[G-22: Fourier frames give stable common tensors in more rows]\label{G-22}
Let $n>m\ge r\ge2$, put $a=n-m$, and let
$\nu\vdash ns$, $s\ge1$, have $r$ positive parts. Suppose
\[
 \nu_r>0,\qquad \nu_1>as.
\]
Set $V=\mathbb C^r$, $\eta=\nu-\nu_r(1,\ldots,1)$,
$e=|\eta|$, and let $v_\eta\in S_\eta V^*$ be the unit
extremal flag vector of weight $-\eta$ used in G-17.
Then the two enhanced cones
\[
 X_Y=\{F^{\otimes s}\otimes v:F\in Y,
                  v\in\mathcal C_\eta\}^{\rm cl}
 \subseteq W=\operatorname{Sym}^s(\operatorname{Sym}^nV)
                         \otimes S_\eta V^*,
\]
where $Y=\mathcal D_{n,r}$ or
$Y=\{L^a q:q\in\mathcal P_{m,r}\}^{\rm cl}$,
have a common stable point for $\mathrm{SL}_r$.
Consequently their invariant coordinate rings have dimensions
\[
 \dim\mathbb C[X_Y]^{\mathrm{SL}_r}
 =\dim Y+r_\eta-(r^2-1),\qquad
 r_\eta=\#\{i<j:\eta_i>\eta_j\},
\]
and their degree-$j$ invariant dimensions equal the respective
multiplicities $m_{j\nu}(Y)$ in original coordinate degree $sj$.
No inequality comparing the dimensions of the two cones is needed.
\end{proposition}

\paragraph{Sources and the additional construction.}
The stability criteria are Losev (2006), Theorem~4, p.~1
(Kempf--Ness), and Mukai (2003), Theorem~7.4, printed p.~212,
with its linear-representation formulation (7.3), p.~213
and proof pp.~218--219 (Hilbert--Mumford). The exact quotient
locators are the same as in G-17. We prove the Fourier-frame
construction and preservation of every one-parameter minimum
weight under multiplication; the moment of a product of linear
forms is not assumed to equal the sum of their separate moments.

\begin{proof}
Write $x_0,\ldots,x_{r-1}$ for an orthonormal basis of $V$ and put
\[
 \mu_0=\frac{\nu_1}{s}-a,\qquad
 \mu_i=\frac{\nu_{i+1}}s\quad(1\le i<r).
\]
Every $\mu_i$ is positive and $\sum_i\mu_i=m$.
Choose a primitive $m$th root $\zeta$, and define the unit vectors
\begin{equation}\label{G22:frame}
 u_j=\sum_{i=0}^{r-1}\sqrt{\frac{\mu_i}{m}}\,
                    \zeta^{ij}x_i,\qquad 0\le j<m.
\end{equation}
The finite Fourier orthogonality identity gives
\begin{equation}\label{G22:framemoment}
 \|u_j\|^2=1,\qquad
 \sum_{j=0}^{m-1}u_ju_j^*=\operatorname{diag}(\mu_0,\ldots,\mu_{r-1}).
\end{equation}
Indeed $\sum_j\zeta^{(i-k)j}$ is $m$ for $i=k$ and zero
otherwise, since $0\le i,k<r\le m$.
Every $r$ of the vectors $u_j$ are independent: after dividing
the rows by their nonzero factors, their determinant is the
Vandermonde determinant on the distinct corresponding powers
of $\zeta$.

First work in the ordered tensor representation
\[
 \widetilde W=V^{\otimes ns}\otimes S_\eta V^*
\]
and consider the nonzero unit decomposable tensor
\[
 \xi=\bigl(x_0^{\otimes a}\otimes
             u_0\otimes\cdots\otimes u_{m-1}\bigr)^{\otimes s}
               \otimes v_\eta.
\]
The moment of a unit tensor product is the sum of the moments
of its factors: expand the infinitesimal action as a sum over
tensor positions and take the product Hermitian inner product.
By \eqref{G22:framemoment}, its diagonal contribution before
removing the trace is
\[
 sa e_0+s\mu-\eta=\nu-\eta=\nu_r(1,\ldots,1).
\]
The off-diagonal entries cancel in the Fourier sum. The
extremal weight vector $v_\eta$ has zero off-diagonal moment,
because distinct torus weights are orthogonal. Hence $\xi$
has zero $\mathrm{SL}_r$ moment map.

Its projective stabilizer is finite. A transformation preserving
the line of an ordered nonzero decomposable tensor preserves
the line of each tensor factor; this follows, for example, by
contracting all other factors with nonzero linear functionals.
If $m\ge r+1$, it consequently fixes $r+1$ full-spark lines
$[u_j]$. The first $r$ are a basis, so the transformation is
diagonal in that basis. Every coordinate of the next vector
in that basis is nonzero, since any $r$ of these $r+1$ vectors
are independent. Fixing its line makes all diagonal entries
equal. Thus the transformation is scalar.
If $m=r$, the $u_j$ are a basis and
\[
 \sum_{j=0}^{m-1}u_j=\sqrt{m\mu_0}\,x_0.
\]
All basis coefficients of $x_0$ are nonzero. Since $a>0$,
the ordered tensor also contains this padding vector and
the same argument again forces a scalar. The projective
stabilizer inside $\mathrm{SL}_r$ is therefore contained in
its finite center. Kempf--Ness now makes the affine orbit
of $\xi$ closed, and its finite stabilizer makes $\xi$ stable.

Let $q=\prod_{j=0}^{m-1}u_j\in\operatorname{Sym}^mV$,
$F=x_0^a q$, and
\[
 w=F^{\otimes s}\otimes v_\eta\in W.
\]
The diagonal $m\times m$ pencil with entries $u_j$ has both
permanent and determinant equal to $q$. Adding $a$ diagonal
entries $x_0$ proves membership of $F$ in both required
degree-$n$ cones. It remains to prove stability after the
symmetrizations; equivariance alone would not prove that.

For any nontrivial one-parameter subgroup
$\lambda:\mathbb G_m\to\mathrm{SL}_r$, diagonalize its action
and denote the smallest occurring integer weight of a nonzero
vector $z$ by $\omega_\lambda(z)$. Smallest weights add under
tensor product. They also add under multiplication of nonzero
polynomials in $\operatorname{Sym}(V)$: their smallest-weight
initial terms have nonzero product because this symmetric
algebra is an integral domain. The same statement applies
to the power $F^{\otimes s}$ in
$\operatorname{Sym}^s(\operatorname{Sym}^nV)$.
It follows that
\begin{equation}\label{G22:minweight}
 \omega_\lambda(w)
 =s\left(a\omega_\lambda(x_0)+
             \sum_j\omega_\lambda(u_j)\right)
                +\omega_\lambda(v_\eta)
 =\omega_\lambda(\xi).
\end{equation}
Both representations have scalar weight
$ns-e=r\nu_r>0$, so they are linear actions of ray type,
with irrelevant set the origin. Mukai's Theorem~7.4 says
that a nonzero vector is stable exactly when its limit
under every nontrivial $\lambda(t)$ as $t\to0$ fails to
exist. In the present weight notation, this means that
every $\omega_\lambda$ is negative. It holds for $\xi$,
and \eqref{G22:minweight} proves it for $w$. Thus $w$ is
the required common stable point.

For completeness the representation-theoretic extraction
uses the same argument as G-17 in rank $r$. The extremal
orbit cone has
$\mathbb C[\mathcal C_\eta]_j=S_{j\eta}V$ and projective
dimension $r_\eta$. The Segre--Veronese map consequently gives
\[
 \dim X_Y=\dim Y+r_\eta,\qquad
 \mathbb C[X_Y]_j=\mathbb C[Y]_{sj}\otimes S_{j\eta}V.
\]
These equalities follow from the span of translates of the
extremal tensor and from multiplication of paired coordinate
functions, as proved in G-17; they do not require normality.
The only degree-$sj$ type whose dual pairs to an invariant
with $S_{j\eta}V$ is
$\lambda=j\eta+t(1,\ldots,1)$.
The size equation $|\lambda|=nsj$ forces $t=j\nu_r$,
so that $\lambda=j\nu$. Exactness of invariants gives the
claimed multiplicity identity. The nonempty stable open
has orbit dimension $r^2-1$, which proves the quotient
dimension formula by the dimension-of-fibers theorem.
\end{proof}

\paragraph{Scope.}
The positive numbers $\mu_i$ need not be rational squares;
this is an existence witness over $\mathbb C$, not a claim
of polynomial bit complexity for a hitting set or for
multiplicity extraction. It works for every stated strict
direction and every $m\ge r$. Removing the scalar-orbit
period of the multiplied tensor requires a further step:
the trivial projective stabilizer of the ordered tensor
does not itself rule out permutations of the factors of $F$.

